%% toc: P15 160 160.1,160.2,160.3,160.4
\chapter{Stirling's Approximation, the Fair-Sign Normal Limit and Three Convolution Details}
\label{sec:p15-w02c-p15-011-stirling-record}

\Cref{sec:p15-de-moivre-laplace-and-stirling} proves Stirling's formula with
Robbins's two-sided remainder and derives the de Moivre--Laplace theorem from
factorial asymptotics. Two parts of the source material behind that chapter
are recorded here. The first is Feller's numerical record of how good the
approximation already is for small \(n\), together with his historical and
bibliographic notes on the formula, its constant and its remainder
\cite[II.9, pp.~52--54]{Feller1968ProbabilityI}. The second is Durrett's
remark that his Theorem~3.1.3 is a special case of the central limit theorem
of his Section~3.4, which is his reason for leaving its details to the reader
\cite[pp.~113, 115]{DurrettPTEAuthor2019}. The third section writes out,
step by step, three arguments that
\cref{sec:p15-convolution-gamma-normal} compresses: the second cumulative
formula for a convolution, the rescaling of the beta integral, and the sum of
a normal variable with a zero-variance normal variable.

\paragraph{Inputs and conclusions.}
The first section takes the closed form
\(A_n=\sqrt{2\pi}\,n^{n+1/2}e^{-n}\) and the remainder bounds of
\cref{thm:p15-stirling-with-remainder}. It concludes with certified values of
the relative errors \(1-A_n/n!\) at \(n=1,2,5,10,100\), the size of the
overestimate made by the upper Robbins member \(A_ne^{1/(12n)}\), and a
comparison of each figure with the one Feller prints. The second section takes
the fair-sign sums \(S_N\) of \cref{def:p15-stirling-models} and the
independent identically distributed central limit theorem
\cref{thm:p15-clt-iid}. It concludes the interval limit of
\cref{thm:p15-stirling-de-moivre-laplace} by a route that does not use
factorials, and it records what that route does not give. The third section
takes product measure and Tonelli's theorem (\cref{thm:p15-tonelli}), the
affine scaling of Lebesgue measure (\cref{thm:p15-affine-lebesgue}) and the
definitions of the convolution law, the beta integral and the normal laws. It
concludes the three identities stated in
\cref{thm:p15-w02c-p15-011-cdf-other-order,thm:p15-w02c-p15-011-beta-rescaling,thm:p15-w02c-p15-011-zero-variance-sum}.

\section{Feller's numerical record and his sources}
\label{sec:p15-w02c-p15-011-feller-record}

\begin{hbconvention}{def:p15-w02c-p15-011-stirling-members}{The Stirling term and the two Robbins members}
For integers \(n\geq1\), \(A_n=\sqrt{2\pi}\,n^{n+1/2}e^{-n}\) is the right
side of Feller's display (9.1), as in
\cref{thm:p15-stirling-worked-boundaries}. By
\cref{thm:p15-stirling-with-remainder}, \(n!=A_ne^{r_n}\) with
\(1/(12n+1)<r_n<1/(12n)\). The \emph{lower} and \emph{upper members} are
\(A_ne^{1/(12n+1)}\) and \(A_ne^{1/(12n)}\); they are the two outer terms of
Feller's double inequality (9.15)
\cite[p.~54]{Feller1968ProbabilityI}. The \emph{relative error} of \(A_n\) is
\(1-A_n/n!\); Feller states it in per cent, that is, multiplied by \(100\).
\end{hbconvention}

\begin{hbtheorem}{thm:p15-w02c-p15-011-stirling-numerics}{proved=annex:p15-w02c-p15-011-stirling-numerics}{Feller's illustrations of Stirling's approximation, certified}
\begin{enumerate}
\item The relative error \(1-A_n/n!\), in per cent, lies in the following
      intervals:
      \[
       \begin{array}{c|c|c}
        n & \text{certified interval (per cent)} & \text{Feller's figure}\\\hline
        1 & (7.786,\,7.787) & 8\\
        2 & (3.92,\,4.17) & 4\\
        5 & (1.62,\,1.67) & 2\\
        10 & (0.823,\,0.834) & 0.8\\
        100 & (0.0832,\,0.0834) & 0.08
       \end{array}
      \]
      Each of Feller's figures is the value of the corresponding interval
      rounded to the number of digits he prints.
\item \(A_1=0.92213\ldots\), \(A_2=1.91900\ldots\) and
      \(A_5=118.0191\ldots\); rounded as Feller prints them, these are
      \(0.9221\), \(1.919\) and \(118.019\). Further
      \(A_{10}=3{,}598{,}695.6\ldots\), against \(10!=3{,}628{,}800\).
\item For every \(n\geq1\) the upper member overestimates \(n!\):
      \[
       0<\frac{A_ne^{1/(12n)}-n!}{n!}
         <e^{1/(12n(12n+1))}-1
         <\frac{0.70}{n^2}\cdot\frac1{100}.
      \]
      That is, the overestimate is less than \(0.70\,n^{-2}\) per cent.
      The exact ratio of the lower member to the upper member is
      \(e^{-1/(12n(12n+1))}\), which exceeds \(1-(12n^2)^{-1}\).
\item The upper member equals \(2.00065\ldots\) at \(n=2\) and
      \(120.00263\ldots\) at \(n=5\).
\end{enumerate}
\index{Stirling's formula!numerical accuracy}
\end{hbtheorem}

Feller uses these figures to separate two notions of error
\cite[p.~52]{Feller1968ProbabilityI}. The absolute difference
\(n!-A_n\) grows without bound, as
\cref{thm:p15-stirling-worked-boundaries} proves, while the percentage error
decreases steadily and is already small for small \(n\). He gives the
approximations \(0.9221\), \(1.919\) and \(118.019\) to \(1!\), \(2!\) and
\(5!=120\), the percentage errors \(8\), \(4\) and \(2\), the approximation
\(3{,}598{,}600\) to \(10!=3{,}628{,}800\) with error \(0.8\) per cent, and
an error of \(0.08\) per cent for \(100!\). After (9.15) he compares the
ratio of its extreme members with \(1-(12n^2)^{-1}\), concludes that the
right-hand member overestimates \(n!\) by an error of order \(n^{-2}\) per
cent, and adds that the true error is much smaller: (9.15) gives
\(2.0007\) for \(n=2\) and \(120.01\) for \(n=5\)
\cite[p.~54]{Feller1968ProbabilityI}.

\paragraph{Reconciling the printed figures.}
Three of the printed figures differ from the exact values in their last
digit, and the theorem keeps the exact values.
\begin{itemize}
\item Feller's \(3{,}598{,}600\) for \(A_{10}\) agrees with the exact
      \(3{,}598{,}695.6\ldots\) to five significant figures; the nearest
      hundred is \(3{,}598{,}700\). His percentage error \(0.8\) is correct
      at the precision printed.
\item At \(n=5\) the upper member is \(120.00263\ldots\).
      \Cref{thm:p15-stirling-worked-boundaries} rounds this to \(120.003\);
      Feller's \(120.01\) exceeds the exact value by about \(0.0074\), and is
      therefore an upward rounding, not the nearest value to two decimals,
      which is \(120.00\). Both are upper bounds for \(5!=120\).
\item Feller's comparison quantity \(1-(12n^2)^{-1}\) is smaller than the
      exact ratio \(e^{-1/(12n(12n+1))}\), so the percentage bound read from
      it, \(100/(12n^2)\) per cent, is valid but larger than item~3's
      \(0.70\,n^{-2}\) per cent. The numerical constant of Feller's printed
      percentage bound is obscured by a mark in the scanned copy consulted
      for this chapter and is not reproduced.
\end{itemize}

\paragraph{Feller's sources for the formula, the constant and the remainder.}
Feller's footnotes record three attributions that the proofs in
\cref{sec:p15-de-moivre-laplace-and-stirling} rely on or replace.
\begin{itemize}
\item The formula is James Stirling's, from his \emph{Methodus
      differentialis} of 1730 \cite[p.~52, n.~14]{Feller1968ProbabilityI}.
\item The proof by the differences \(d_n-d_{n+1}\), and the double
      inequality (9.14) behind the Robbins bounds of
      \cref{thm:p15-stirling-with-remainder}, are due to H.~E. Robbins,
      \emph{Amer. Math. Monthly}, vol.~62 (1955), pp.~26--29, as cited by
      Feller \cite[p.~53, n.~15]{Feller1968ProbabilityI}.
\item The usual proof that the constant \(e^C\) equals \(\sqrt{2\pi}\)
      relies on Wallis's formula. Feller instead determines the constant in
      VII.2 by a probabilistic argument, and for a simple direct proof he
      refers to his own note in \emph{Amer. Math. Monthly} (1967)
      \cite[p.~53, n.~16]{Feller1968ProbabilityI}. The handbook's
      determination is Feller's VII.2 argument, with an independent second
      determination from the cosine-integral central coefficient, both in
      \cref{thm:p15-stirling-with-remainder}.
\end{itemize}
The refinement \(d_n=C+(12n)^{-1}-(360n^3)^{-1}+\cdots\) mentioned in
Feller's footnote~17 \cite[p.~54, n.~17]{Feller1968ProbabilityI} is the
explicit fifth-order remainder bound proved in
\cref{thm:p15-stirling-with-remainder}.

\begin{hbwarning}{def:p15-w02c-p15-011-relative-not-absolute}{Small percentage error is not small absolute error}
Item~1 shows the relative error falling below one per cent by \(n=10\), yet
the absolute error \(10!-A_{10}=30{,}104.3\ldots\) already exceeds
\(30{,}000\), and \(n!-A_n\to\infty\) by
\cref{thm:p15-stirling-worked-boundaries}. Stirling's formula controls only
the ratio \(n!/A_n\); it cannot by itself decide a comparison between \(n!\)
and another integer near it. The upper member in item~3 is an upper bound
for every \(n\), and at \(n=5\) it lies within \(0.003\) of \(120\). Its
distance from \(n!\) is controlled only through the ratio bound of item~3,
that is, as a fraction of \(n!\).
\end{hbwarning}

\section{The fair-sign limit as a case of the central limit theorem}
\label{sec:p15-w02c-p15-011-dml-from-clt}

Durrett opens his chapter on central limit theorems by treating special cases
that can be handled by elementary computation, and after stating his
Theorem~3.1.3 he notes that it is a special case of the central limit theorem
of his Section~3.4, which is why he leaves its remaining details to the
reader \cite[pp.~113, 115]{DurrettPTEAuthor2019}. Those details are written
out by the factorial route in \cref{thm:p15-stirling-de-moivre-laplace}. The
corollary below proves the same interval limit from \cref{thm:p15-clt-iid},
whose proof uses characteristic functions and no factorial estimate.

\begin{hbcorollary}{thm:p15-w02c-p15-011-dml-from-clt}{proved=annex:p15-w02c-p15-011-dml-from-clt}{Durrett's Theorem~3.1.3 from the general central limit theorem}
Let \(\varepsilon_1,\varepsilon_2,\ldots\) be independent fair signs and
\(S_N=\sum_{i=1}^N\varepsilon_i\), as in \cref{def:p15-stirling-models}.
Then \(\mathbb E\varepsilon_1=0\), \(\operatorname{Var}(\varepsilon_1)=1\),
and \(S_N/\sqrt N\Rightarrow\chi\), where \(\chi\) has the standard normal
law. Consequently, for all \(-\infty\leq a\leq b\leq+\infty\),
\[
 \mathbb P(a\leq S_N/\sqrt N\leq b)\longrightarrow\Phi(b)-\Phi(a)
 =\int_a^b\phi(x)\,dx,
\]
with \(\Phi(-\infty)=0\), \(\Phi(+\infty)=1\), and the same limit for each
choice of strict or non-strict endpoint inequalities.
\index{de Moivre--Laplace theorem!as a central limit theorem}
\end{hbcorollary}

\begin{hbexample}{def:p15-w02c-p15-011-what-the-clt-misses}{What the general theorem does not give}
The corollary with \(a=b=0\) gives \(\mathbb P(S_{2n}=0)\to0\), but no rate.
The factorial route gives the rate: by the exact central coefficient and
\cref{thm:p15-stirling-with-remainder},
\[
 \mathbb P(S_{2n}=0)=\binom{2n}n4^{-n}
  =\frac{(2n)!}{(n!)^2}4^{-n}
  =\frac{e^{r_{2n}-2r_n}}{\sqrt{\pi n}},
\]
so \(\sqrt{\pi n}\,\mathbb P(S_{2n}=0)\to1\). At \(n=5\),
\(\binom{10}5/1024=252/1024=0.24609\ldots\), while
\(1/\sqrt{5\pi}=0.25231\ldots\). The local and rate statements of
\cref{thm:p15-stirling-local-rademacher} are therefore not consequences of
\cref{thm:p15-clt-iid}; the central limit theorem controls probabilities of
intervals of length of order \(\sqrt N\), not of single lattice points.
\end{hbexample}

\paragraph{Relationships.}
The two proofs of the interval limit are independent: the factorial proof
uses \cref{thm:p15-stirling-local-rademacher}, and the characteristic-function
proof uses \cref{thm:p15-clt-iid}. The fixed-parameter binomial limit of
\cref{thm:p15-stirling-general-binomial} follows from \cref{thm:p15-clt-iid}
in the same way: the success indicators are independent with mean \(p\)
and variance \(pq>0\), so \((B_N-Np)/\sqrt{Npq}\Rightarrow\chi\). Two of the failures recorded in
\cref{thm:p15-stirling-worked-boundaries}, dependent signs and success
probabilities changing with \(N\), violate the independent identically
distributed hypothesis of \cref{thm:p15-clt-iid} as well. The third, the loss
of relative accuracy at the edge of the support, concerns a uniform relative
statement that neither theorem makes.

\section{Three convolution details written out}
\label{sec:p15-w02c-p15-011-convolution-details}

Durrett states the law of an independent sum with the integration taken
against the second variable only,
\(\mathbb P(X+Y\leq z)=\int F(z-y)\,dG(y)\) in his Theorem~2.1.15
\cite[pp.~49--50]{DurrettPTEAuthor2019}.
\Cref{thm:p15-convolution-law} also states the other order. In the proof of
his Theorem~2.1.18 he obtains the rescaled beta integral by the change of
variables \(y=xu\), \(dy=x\,du\), display (2.1.1)
\cite[p.~51]{DurrettPTEAuthor2019}, and his Theorem~2.1.20 treats normal
summands with positive variances only \cite[p.~52]{DurrettPTEAuthor2019}.
The zero-variance law \(\operatorname{normal}(\nu,0)=\delta_\nu\) is the
declared extension of \cref{def:p15-normal-law}. The three results below
isolate the steps of \cref{thm:p15-convolution-law},
\cref{thm:p15-beta-gamma-convolution} and \cref{thm:p15-normal-sums} that
concern these points, and their annex proofs give every step. As in those
two theorems, \(\mu\) and \(\nu\) denote laws in the first lemma and normal
means in the third.

\begin{hblemma}{thm:p15-w02c-p15-011-cdf-other-order}{proved=annex:p15-w02c-p15-011-cdf-other-order}{The cumulative formula integrated against the first law}
Let \(\mu,\nu\) be Borel probability measures on \(\mathbb R\) with
cumulative functions \(F,G\), and let \(z\in\mathbb R\). Put
\(H_z=\{(x,y)\in\mathbb R^2:x+y\leq z\}\). Then \(H_z\) belongs to the
product sigma-algebra \(\mathcal B(\mathbb R)\times\mathcal B(\mathbb R)\);
for each fixed \(x\),
\[
 \int_{\mathbb R}\mathbf1_{H_z}(x,y)\,d\nu(y)=\nu((-\infty,z-x])=G(z-x);
\]
the function \(x\mapsto G(z-x)\) is Borel measurable; and
\[
 (F*G)(z)=(\mu\times\nu)(H_z)=\int_{\mathbb R}G(z-x)\,d\mu(x)
         =\int_{\mathbb R}G(z-x)\,dF(x),
\]
where, as in \cref{def:p15-convolution-law}, \(dF(x)\) means integration
against \(\mu\).
\end{hblemma}

As a check, take \(\nu=\delta_c\). Then \(G(z-x)=1\) when \(c\leq z-x\)
and \(0\) otherwise, so the lemma gives
\((F*G)(z)=\mu((-\infty,z-c])=F(z-c)\), which agrees with the translation
rule \((\mu*\delta_c)(A)=\mu(A-c)\) of \cref{thm:p15-convolution-law} at
\(A=(-\infty,z]\). With the roles of the two laws exchanged, the same
argument gives Durrett's order \(\int F(z-y)\,dG(y)\); the annex proof of
\cref{thm:p15-convolution-law} writes that order out.

\begin{hblemma}{thm:p15-w02c-p15-011-beta-rescaling}{proved=annex:p15-w02c-p15-011-beta-rescaling}{Rescaling the beta integral to the interval \((0,z)\)}
Let \(\alpha,\beta>0\) and \(z>0\). With the zero-endpoint convention of
\cref{def:p15-beta-integral},
\[
 \int_0^z(z-y)^{\alpha-1}y^{\beta-1}\,d\lambda(y)
  =z^{\alpha+\beta-1}\int_0^1(1-u)^{\alpha-1}u^{\beta-1}\,d\lambda(u)
  =z^{\alpha+\beta-1}B(\beta,\alpha)
\]
as an equality of nonnegative extended real numbers. Here \(\lambda\) is
Lebesgue measure.
\end{hblemma}

By \cref{thm:p15-beta-gamma-convolution}, \(B(\beta,\alpha)\) is finite and
positive, so both sides are finite and positive. At \(\alpha=\beta=1\) both
sides equal \(z\), since \(B(1,1)=1\). The factor \(z\) contributed by the
substitution is the factor that Durrett writes as \(dy=x\,du\); without it the
right side would carry \(z^{\alpha+\beta-2}\), which at \(\alpha=\beta=1\)
gives \(1\) instead of \(z\), wrong for every \(z\neq1\).

\begin{hblemma}{thm:p15-w02c-p15-011-zero-variance-sum}{proved=annex:p15-w02c-p15-011-zero-variance-sum}{A zero-variance normal summand}
Let \(X\) have law \(\operatorname{normal}(\mu,a)\) with \(a\geq0\), and let
\(Y\) have law \(\operatorname{normal}(\nu,0)=\delta_\nu\), on the same
probability space. Then \(Y=\nu\) almost surely, \(X+Y=X+\nu\) almost surely,
and \(X+Y\) has law \(\operatorname{normal}(\mu+\nu,a)\). This is the case
\(b=0\) of the sum formula in \cref{thm:p15-normal-sums}. The conclusion
does not require independence of \(X\) and \(Y\), and in any case \(Y\) is
independent of \(X\).
\end{hblemma}

The independence hypothesis of \cref{thm:p15-normal-sums} is therefore
automatic when one summand has variance zero. The contrast with positive variances is
\cref{thm:p15-convolution-calculations-boundaries}, item~4, where two standard
normal summands that are not independent have sum \(\delta_0\) rather than
\(\operatorname{normal}(0,2)\).
